\section{Deep Representation Learning}\label{sec:bg:dl}

Many \gls{dl} pipelines replace hand-crafted features with a parametric map that is trained end-to-end from data.
The central object is an \emph{encoder} \(\Enc \colon \setX \to \R^{d}\), implemented in practice as a stack of nonlinear layers (for example, convolutions over a grid, recurrence over time, or attention over a sequence) followed by a pooling or readout head that yields a fixed-dimensional embedding \(\vect{z} = \Enc(\vect{x})\).
The design question is which training signal makes \(\vect{z}\) useful for the downstream task when labels are scarce or absent.

In supervised learning, \(\Enc\) is trained together with a classifier so that \(\vect{z}\) separates classes under a cross-entropy or margin loss.
When cluster structure is the target but assignments are unknown, the same encoder can instead be trained with objectives that encourage smoothness, invariance to nuisance transformations, or agreement with a clustering hypothesis in embedding space.
\emph{Self-supervised} learning occupies an intermediate regime: the model predicts surrogate targets constructed from the input itself (for example, masked tokens, relative positions, or transformed views), so that \(\Enc\) learns semantics without human labels \parencite{chen2020simclr}.

\emph{Contrastive} objectives implement instance discrimination by treating two augmentations of the same example as a positive pair and other examples in a minibatch (or a memory bank) as negatives.
Let \(\vect{z}_i\) and \(\vect{z}_{i}^{+}\) denote normalized embeddings of two views of sample \(i\), let \(\mathrm{sim}(\cdot,\cdot)\) denote cosine similarity, and let \(\tau>0\) be a temperature.
A standard \emph{InfoNCE}-style loss takes the form
\begin{equation}
\label{eq:bg:infonce}
\Loss_{\mathrm{NCE}}(i)
=
-\log
\frac{\exp\bigl(\mathrm{sim}(\vect{z}_i,\vect{z}_{i}^{+})/\tau\bigr)}
{\exp\bigl(\mathrm{sim}(\vect{z}_i,\vect{z}_{i}^{+})/\tau\bigr)
 + \sum_{j \neq i} \exp\bigl(\mathrm{sim}(\vect{z}_i,\vect{z}_{j}^{+})/\tau\bigr)},
\end{equation}
which lower-bounds mutual information between representations of different views under suitable generative assumptions \parencite{oord2018representation}.
Large-scale implementations such as SimCLR combine strong augmentations with wide embeddings and many negatives drawn from the same minibatch \parencite{chen2020simclr}.
\cref{eq:bg:infonce} should be read as a template: the positive set can be enlarged, negatives can be drawn from a queue, and the similarity can be implemented with a learned projection head without changing the underlying contrastive principle.

Pretext tasks offer an alternative route: the network solves an auxiliary problem (predicting rotations, solving jigsaw permutations, inpainting masked regions) whose solution is incidental, but the representation \(\vect{z}\) becomes predictive of semantically stable factors.
Such objectives can be combined with clustering when the end goal is partition discovery rather than linear probe accuracy on a fixed label set.

\Gls{dec} couples an encoder with a soft assignment to a small set of cluster centroids in embedding space and trains the model by minimizing a Kullback--Leibler divergence between the soft assignments and an auxiliary target distribution that sharpens confident predictions \parencite{xie2016unsupervised}.
SwAV replaces pairwise instance contrast with \emph{prototype} classification: views are mapped to cluster codes, and the loss enforces prediction consistency across augmentations while preventing collapse through swapped assignments \parencite{caron2020swav}.
Both families treat the encoder as the shared backbone through which raw inputs are mapped to a geometry where classical clusterers or learned assignment heads operate.

Sequence encoders based on self-attention, including the transformer blocks reviewed in \cref{sec:bg:transformer}, instantiate \(\Enc\) when \(\vect{x}\) is an ordered set of tokens (for example, a window of \glspl{pdw}).
The Method chapter builds on this stack: a contrastive or clustering-friendly loss shapes the embedding space, while the attention-based encoder models temporal context within each window.
